\chapter{Vertex algebras and the lattice laboratory}
\label{ch:vertex_algebras}

The preceding chapter described fields through their correlation
functions and conformal transformations. A vertex algebra packages
their operator products algebraically. A lattice vertex algebra is a
particularly concrete example: its states are oscillator monomials
attached to lattice momenta, and its products follow from one
exponential formula. This is the setting in which the Haagerup
construction becomes explicit.

We begin with the space of chiral states and explain how a field acts
on that space. The distinction between a state, a field, and one
coefficient of a field is essential: they are three different kinds of
mathematical object. Once their types and notation are clear, the
vertex-algebra axioms express the vacuum, translation, and locality
properties familiar from CFT. We then build a lattice example in which
the states and the operators can both be written explicitly.

\section{The state--field map and the mode dictionary}
\index{vertex operator algebra}\index{modes!vertex-algebra convention}

\subsection{States, formal fields, and coefficient operators}
\index{state--field map}\index{formal Laurent series}

Let $V$ be a complex vector space of chiral states. For the CFT-type
examples below, it will be the algebraic direct sum of the finite
energy spaces in the vacuum sector. An element of this direct sum has
components in only finitely many energy spaces; a Hilbert-space completion
is additional structure, introduced when we discuss unitarity.
Write $\End_{\C}(V)$ for the space of linear maps $V\to V$.
Its elements are operators on states.

For now $z$ is a \emph{formal variable}. A series in $z$ means a
specified list of coefficients, with algebraic rules for extracting,
differentiating, and comparing them. It does not yet mean a function
evaluated at a nonzero complex number. In particular, no convergence
or boundedness of operators is assumed in the following definitions.
We write
\[
 V[[z]]=\left\{\sum_{j\geq0}v_jz^j:v_j\in V\right\},
 \qquad V((z))=\bigcup_{N\geq0}z^{-N}V[[z]].
\]
Thus $V((z))$ consists of vector-valued Laurent series with only
finitely many negative powers, but possibly infinitely many positive
powers. Neither space is the original state space $V$.

A \emph{field on $V$} is an operator-valued formal series
\[
 A(z)=\sum_{n\in\Z}A_nz^{-n-1},\qquad A_n\in\End_{\C}(V),
\]
such that, for every $b\in V$, one has $A_nb=0$ for all sufficiently
large $n$. This is \emph{lower truncation}: after acting on $b$, the
series $A(z)b$ belongs to $V((z))$. The bound on $n$ can depend on
$b$; there need not be a single bound for the operators $A_n$ on all
of $V$. The indexing $-n-1$ is conventional. Large positive $n$
correspond to large negative powers of $z$.

Denote the vector space of such fields by $\mathcal F(V)$. The
\emph{state--field map} is a linear map
\begin{equation}
 Y:V\longrightarrow\mathcal F(V),\qquad a\longmapsto Y(a,z).
 \label{eq:voa-state-field-type}
\end{equation}
It is part of the vertex-algebra structure, just as multiplication
is part of an associative algebra's structure. The vector space
$V$ alone does not determine it. Equivalently, regarding a field
as its action on an input state, one can write
$Y:V\to\Hom_{\C}(V,V((z)))$.
The common notation
$Y(a,z)\in\End_{\C}(V)[[z,z^{-1}]]$ uses the vector space of
doubly infinite formal distributions. It allows powers in both
directions; the separate lower-truncation condition is still required.
It does not authorize arbitrary products of two series at the same
variable. The normal product below gives the relevant prescription.

For a given state $a$, denote the coefficient \emph{operator} of
$z^{-n-1}$ in its field by $a_{(n)}$, and call it the $n$th
\emph{mode} of $a$. Thus
\begin{equation}
 Y(a,z)=\sum_{n\in\Z}a_{(n)}z^{-n-1},
 \qquad a_{(n)}:V\longrightarrow V.
 \label{eq:voa-mode-definition}
\end{equation}
The symbol $a_{(n)}$ is neither a scalar coordinate of $a$ nor
another state. Applying it to a state produces a state:
$a_{(n)}b\in V$. Linearity of $Y$ implies that the product
$(a,b)\mapsto a_{(n)}b$ is bilinear for every fixed integer $n$.
There is a whole family of these products, indexed by $n$.

The types can be kept in view as follows:
\begin{center}
\begin{tabular}{ll}
\toprule
Object & Where it lives or what it does\\
\midrule
$a,b$ & States in $V$\\
$Y$ & Linear map from states to fields\\
$Y(a,z)$ & Formal series of operators on $V$\\
$a_{(n)}$ & One linear operator $V\to V$\\
$Y(a,z)b$ & Vector-valued series in $V((z))$\\
$a_{(n)}b$ & One state in $V$\\
\bottomrule
\end{tabular}
\end{center}

For example, writing out a field's action makes the index shift visible:
\[
 Y(a,z)b=\cdots+(a_{(1)}b)z^{-2}+(a_{(0)}b)z^{-1}
                  +a_{(-1)}b+(a_{(-2)}b)z+\cdots.
\]
The negative-power tail on the left is finite for this particular
$b$. The constant coefficient is $a_{(-1)}b$, whereas the residue
is $a_{(0)}b$. In general,
\[
 a_{(n)}b=\operatorname{Res}_z\bigl(z^nY(a,z)b\bigr),
\]
where the \emph{formal residue} means the coefficient of $z^{-1}$.
It can be computed without choosing a contour or proving convergence.

\subsection{The vacuum, translation, and locality axioms}

To make these fields into a vertex algebra, specify a distinguished
nonzero state $\one\in V$, the \emph{vacuum}, and a linear
operator $\mathsf T:V\to V$, the \emph{translation operator}.
The vacuum and translation identities are
\begin{align}
 Y(\one,z)&=\mathrm{id},&
 Y(a,z)\one&=e^{z\mathsf T}a,\\
 \mathsf T\one&=0,&
 [\mathsf T,Y(a,z)]&=Y(\mathsf Ta,z)=\partial_zY(a,z).
 \label{eq:voa-vacuum-translation}
\end{align}
Here $\mathrm{id}$ is the identity \emph{operator} on $V$, while
$\one$ is a \emph{state}. The exponential means the formal series
$e^{z\mathsf T}a=\sum_{r\geq0}z^r\mathsf T^ra/r!$.
The commutator $[A,B]=AB-BA$ uses composition of linear operators;
the derivative on the right is taken coefficient by coefficient.
Translation therefore differentiates a field's position variable,
as it does in the physical state--operator correspondence.

\begin{example}[Recovering a state from its field]
\label{ex:voa-vacuum-modes}
Compare coefficients in $Y(a,z)\one=e^{z\mathsf T}a$. There
are no negative powers on the right, so
\[
 a_{(n)}\one=0\quad(n\geq0),\qquad
 a_{(-r-1)}\one=\frac{\mathsf T^ra}{r!}\quad(r\geq0).
\]
In particular $a_{(-1)}\one=a$ and $a_{(-2)}\one=\mathsf Ta$.
Thus the field determines its state: act on the vacuum and take
the constant coefficient. If $Y(a,z)=0$, this procedure gives
$a=0$, so $Y$ is injective. It need not map onto every formal
field on $V$.

For the vacuum field itself, $Y(\one,z)=\mathrm{id}$ has only
one nonzero mode:
\[
 \one_{(-1)}=\mathrm{id},\qquad \one_{(n)}=0\quad(n\ne-1).
\]
This is why a vacuum state appearing as an OPE coefficient becomes
an identity operator in the corresponding field formula.
\end{example}

Finally, locality says that for any $a,b$ there is an integer $N\geq0$
such that
\[
 (z-w)^N[Y(a,z),Y(b,w)]=0
\]
as an operator-valued formal distribution in two variables.
The same $N$ works on every input state for this pair of fields.
Different expansions in $z/w$ and $w/z$ encode the two radial
orderings; multiplication by $(z-w)^N$ removes the singularity
that prevents these ordered expansions from agreeing. For example,
the two formal expansions of a simple pole are
\[
 \left.\frac1{z-w}\right|_{|z|>|w|}
   =\sum_{r\geq0}w^rz^{-r-1},\qquad
 \left.\frac1{z-w}\right|_{|w|>|z|}
   =-\sum_{r\geq0}z^rw^{-r-1}.
\]
The domain labels recall their analytic geometric-series meanings;
here they specify which expansion to use. The two series differ,
but multiplying either by $z-w$ gives one. Locality is therefore
compatible with a nonzero commutator of the unmultiplied fields.

These are the vacuum, translation, and locality requirements for
the fields assigned by $Y$; an arbitrary assignment of formal
fields to states need not satisfy them.

\subsection{From state products to operator products}

The connection to the OPE notation of Chapter~\ref{ch:cft} is
\begin{equation}
 Y(a,z)Y(b,w)\sim\sum_{n\geq0}
       \frac{Y(a_{(n)}b,w)}{(z-w)^{n+1}}.
 \label{eq:voa-state-product-ope}
\end{equation}
The symbol $\sim$ retains the singular terms at $z=w$ in the
chosen radial ordering. Each numerator on the right is a
\emph{field}; it is obtained by first forming the \emph{state}
$a_{(n)}b$ and then applying $Y$. The sum of singular terms is
finite by lower truncation. The regular terms also have meaning:
the first one is the field of $a_{(-1)}b$, the normal product
discussed in the next section.

The combination of these axioms implies the Jacobi, or Borcherds,
identity, which expresses compatibility of composing fields with
taking their state products. One useful consequence is the mode
commutator formula
\begin{equation}
 [a_{(m)},b_{(n)}]
 =\sum_{j\ge0}\binom mj
          (a_{(j)}b)_{(m+n-j)}.
 \label{eq:voa-mode-commutator}
\end{equation}
Here $(a_{(j)}b)_{(m+n-j)}$ means a mode of the field assigned
to the state $a_{(j)}b$; it does not mean a product of two mode
operators. For $j\geq0$ the binomial coefficient is
$\binom mj=m(m-1)\cdots(m-j+1)/j!$, with $\binom m0=1$;
this definition also applies to negative $m$.
The sum is finite because the singular products $a_{(j)}b$ vanish
for large $j$. Thus knowledge of the singular OPE determines the
commutators of all modes. It does not determine an arbitrary abstract
vertex algebra unless the generators and their relations are also
controlled.

\subsection{Conformal grading and the physics mode convention}

A \emph{conformal vector} is a distinguished state $\omega\in V$
whose field is the stress tensor. Write that field as
\[
 Y(\omega,z)=T(z)=\sum_{m\in\Z}L_mz^{-m-2}.
\]
Comparing exponents with \eqref{eq:voa-mode-definition} defines
$L_m=\omega_{(m+1)}$. To obtain a vertex operator algebra, or
VOA, require these operators to satisfy the Virasoro relations,
with $L_{-1}=\mathsf T$, and require $L_0$ to give an integer
grading bounded below with finite-dimensional graded pieces.
We mainly use VOAs of CFT type:
\[
 V=\bigoplus_{n\ge0}V_n,\qquad V_0=\C\one.
\]
Here $V_n=\{v\in V:L_0v=nv\}$, and a state in a single $V_n$
is \emph{homogeneous}, of weight $\wt(v)=n$. In particular
$\omega$ has weight two. The grading compatibility identity
\[
 [L_0,Y(a,z)]=Y(L_0a,z)+z\partial_zY(a,z)
\]
gives, for homogeneous $a$,
$[L_0,a_{(n)}]=(\wt(a)-n-1)a_{(n)}$.
Acting on a homogeneous state $b$ proves the weight rule
\begin{equation}
 \wt(a_{(n)}b)=\wt(a)+\wt(b)-n-1.
 \label{eq:voa-mode-weight}
\end{equation}
The displayed weight is the weight of a nonzero output; if the
corresponding graded space is zero, the product must vanish.
This rule provides a quick consistency check on every claimed mode
identity. Two weight-two states have a weight-two first product
$a_{(1)}b$ and a weight-three zeroth product $a_{(0)}b$.

Physics often writes a field of weight $h$ as
$\phi(z)=\sum_m\phi_m z^{-m-h}$. The dictionaries agree when
$\phi_m=a_{(m+h-1)}$. In particular, a current of weight one has
$J_m=J_{(m)}$, while a stress of weight two has
$L_m=\omega_{(m+1)}$. For a weight-$h$ state, the physics zero
mode is $a_{(h-1)}$: it preserves conformal weight. The
vertex-algebra zeroth product instead uses $a_{(0)}$, which
changes weight by $h-1$. In particular, the stress has
$\omega_{(0)}=L_{-1}$ and $\omega_{(1)}=L_0$.
Keeping the convention explicit prevents a translation operator
from being mistaken for an energy operator.

\begin{example}[A current commutator from one OPE coefficient]
Let $J\in V_1$ be a current state and write $J_m=J_{(m)}$.
Suppose $J_{(1)}J=\one$ and $J_{(j)}J=0$ for $j\ne1$, $j\ge0$.
Equation~\eqref{eq:voa-mode-commutator} gives
\[
 [J_m,J_n]=m\one_{(m+n-1)}=m\delta_{m+n,0}\,\mathrm{id},
\]
because $Y(\one,z)=\mathrm{id}$ has only the mode
$\one_{(-1)}=\mathrm{id}$. This is the Heisenberg algebra.
For clarity, $J$ here is a weight-one \emph{state},
$J(z)=Y(J,z)$ is its field, and $J_m=J_{(m)}$ is an operator.
The vacuum identities give $J_{-1}\one=J$ and $J_1\one=0$,
so the commutator also yields the state calculation
\[
 J_1J=J_1J_{-1}\one
     =[J_1,J_{-1}]\one+J_{-1}J_1\one=\one.
\]
The OPE coefficient $J_{(1)}J=\one$ and the central term
$[J_1,J_{-1}]=\mathrm{id}$ express the same structure at the
state and operator levels. More generally,
$J_mJ_{-m}\one=m\one$ for every $m>0$. Thus $Y(J,z)$ acting
on $J_{-m}\one$ has a pole of order $m+1$. This gives a concrete
example of lower truncation with no uniform pole bound over all
input states.
\end{example}

\section{Normal products, derivatives, and strong generation}
\index{normal product}\index{strong generation}\index{skew symmetry}

The constant term in the state series $Y(a,z)b$ is $a_{(-1)}b$.
Its corresponding field is the \emph{normal product}:
\[
 Y(a_{(-1)}b,z)=:Y(a,z)Y(b,z):.
\]
To define the colons, split the first field into
\[
 Y(a,z)_-=\sum_{n<0}a_{(n)}z^{-n-1},\qquad
 Y(a,z)_+=\sum_{n\ge0}a_{(n)}z^{-n-1},
\]
and put
\[
 :Y(a,z)Y(b,z):
   =Y(a,z)_-Y(b,z)+Y(b,z)Y(a,z)_+.
\]
This places the creation part on the left and the part that
annihilates the vacuum on the right. Individual modes need not
annihilate every other state. Lower truncation makes each
coefficient of the displayed expression well defined on an input
state, even though the fields themselves are infinite series.
The equality with $Y(a_{(-1)}b,z)$ is a vertex-algebra identity,
linking this operator prescription to the state product.
In particular,
$\mathsf Ta=a_{(-2)}\one$. Differentiating the mode expansion gives
\begin{equation}
 (\mathsf Ta)_{(n)}=-n\,a_{(n-1)}.
 \label{eq:voa-derivative-mode}
\end{equation}
Repeated negative modes can therefore be expressed using derivatives
and normal products. They are not generally ordinary associative
commutative products: changing parentheses or reversing factors adds
terms governed by the singular OPE.

Skew symmetry is the identity
\begin{equation}
 Y(a,z)b=e^{z\mathsf T}Y(b,-z)a.
 \label{eq:voa-skew}
\end{equation}
For example, extracting the coefficient of $z^{-n-1}$ gives
\[
 a_{(n)}b=\sum_{j\ge0}\frac{(-1)^{n+j+1}}{j!}
          \mathsf T^j(b_{(n+j)}a).
\]
The extra derivatives are precisely why a naive commutative algebra
of OPE coefficients loses information.

\begin{definition}
A set $S\subset V$ \emph{generates} $V$ as a vertex algebra if all
iterated modes of its elements, applied to the vacuum, span $V$.
It \emph{strongly generates} $V$ if derivatives and iterated
normal products of its elements span $V$.
\end{definition}

Strong generation is the more restrictive statement. A set can
generate a new low-weight state through singular products of
higher-weight composites. Knowing a few low graded pieces does not
exclude this possibility. For the Haagerup construction, proving that
six chosen currents remain all of $U_1$ requires an all-weight
argument. The finite-closure chapter explains how a controlled set
of identities supplies that argument.

\begin{example}[Why a weight-one conclusion follows from strong generation]
Suppose a CFT-type VOA is strongly generated by six weight-one
currents and fields of weights at least two. Every derivative raises
weight, and every normal product adds weights. Hence a weight-one
state in their normal span can only be a linear combination of the
six currents. Once strong generation is proved, the claim
$\dim V_1\le6$ is immediate. Without strong generation, the same
count is not justified.
\end{example}

\section{The Heisenberg Fock space}
\index{Heisenberg algebra}\index{Fock space}

Let $\mathfrak h$ be an $r$-dimensional complex vector space with a
nondegenerate symmetric form $(\ ,\ )$, obtained by complexifying a
real Euclidean space. For each $a\in\mathfrak h$ and $n\in\Z$,
introduce a mode operator $a(n)$, linear in $a$. These oscillators obey
\[
 [a(m),b(n)]=m(a,b)\delta_{m+n,0}\,\mathrm{id}.
\]
The vacuum representation is generated from a vector $\one$ killed
by $a(n)$ for $n\ge0$. Negative modes commute with one another
and act freely as creation operators. Thus its state space is the
symmetric-algebra Fock space
\[
 M(1)=S(\mathfrak h\otimes t^{-1}\C[t^{-1}]).
\]
Here $S(W)$ means the polynomial, or symmetric, algebra on a vector
space $W$. The letter $t$ labels oscillator modes and is distinct
from the field variable $z$. The symbol $a\otimes t^{-n}$, $n>0$,
labels a creation mode $a(-n)$; a monomial represents a state
$a_1(-n_1)\cdots a_s(-n_s)\one$.
Positive modes act by commuting through these creation modes and
then killing the vacuum. For example,
$a(1)b(-1)\one=(a,b)\one$.

For $a\in\mathfrak h$, the associated weight-one state is
$a(-1)\one$, and its field is
\[
 Y(a(-1)\one,z)=\sum_{n\in\Z}a(n)z^{-n-1}.
\]
Thus $a(n)$ abbreviates $(a(-1)\one)_{(n)}$; the vector
$a\in\mathfrak h$, the state $a(-1)\one\in M(1)$, and its
mode operators have now been related explicitly.
An oscillator $a(-n)$ contributes weight $n$. For an orthonormal
basis $u_1,\ldots,u_r$, the conformal vector is
\[
 \omega=\frac12\sum_{j=1}^r u_j(-1)^2\one,
\]
and the central charge is $r$.

The character is the graded trace
$\ch_{M(1)}(q)=\Tr_{M(1)}q^{L_0-r/24}$.
Its product formula follows from one geometric series for every
oscillator species and frequency:
\begin{equation}
 \ch_{M(1)}(q)
 =q^{-r/24}\prod_{n\ge1}(1-q^n)^{-r}
 =\eta(\tau)^{-r}.
 \label{eq:voa-oscillator-character}
\end{equation}
Here $q=e^{2\pi i\tau}$, with $\operatorname{Im}\tau>0$ for an
analytic character, and
$\eta(\tau)=q^{1/24}\prod_{n\ge1}(1-q^n)$ is the Dedekind eta
function. One can also read the product coefficient by coefficient
as a formal graded count. This $q$ records energy; it is not the
insertion coordinate $z$ of a field.
The Heisenberg algebra alone has a Fock module for every
momentum in $\mathfrak h$; its natural representation theory is not
finite. A positive even lattice supplies a mutually local discrete
set of exponential fields and changes this situation substantially.

\section{Even lattices and cocycles}
\index{lattice vertex algebra}\index{cocycle}

A lattice $L$ is a free abelian group of finite rank with a
positive-definite integral form: $(\alpha,\beta)\in\Z$ for lattice
vectors $\alpha,\beta$. Thus, after choosing a lattice basis, each
vector is an integer linear combination of basis vectors. Write
$\alpha^2=(\alpha,\alpha)$ and set
$\mathfrak h=L\otimes_{\Z}\C$, so the Heisenberg space is the
complexified lattice space.
The lattice is \emph{even} when
$(\alpha,\alpha)\in2\Z$ for every $\alpha\in L$.
We will see how integrality and evenness enter the field construction
after defining its states.

\subsection{Momentum labels, states, and their fields}

The Heisenberg vacuum space $M(1)$ describes oscillator excitations
above a state of zero momentum. To include the momenta in $L$, we
introduce one new \emph{basis symbol} $e^\alpha$ for each
$\alpha\in L$ and form the vector space
\[
 \C_\epsilon[L]=\bigoplus_{\alpha\in L}\C e^\alpha.
\]
The subscript $\epsilon$ anticipates a multiplication to be defined
below; it does not affect this vector-space definition. An element
is a finite sum $\sum_\alpha c_\alpha e^\alpha$ with complex
coefficients. The vector $\alpha$ belongs to the lattice, whereas
$e^\alpha$ belongs to this new space of momentum labels. Here the
notation names a basis vector; no power-series exponential of
$\alpha$ is being taken.

Attach a copy of the oscillator space to each momentum label. The
resulting vector space of the lattice VOA is
\begin{equation}
 \begin{split}
 V_L&=M(1)\otimes\C_\epsilon[L]\\
    &=\bigoplus_{\alpha\in L}\bigl(M(1)\otimes\C e^\alpha\bigr).
 \end{split}
 \label{eq:voa-lattice-space}
\end{equation}
The ground state in the momentum-$\alpha$ summand is
$|\alpha\rangle=\one\otimes e^\alpha$, with no oscillator
excitations. We conventionally abbreviate this state to $e^\alpha$
as well. The vacuum of $V_L$ is $\one\otimes e^0$, also denoted
$\one$. Thus the symbol $e^0$ labels the zero-momentum basis vector;
its tensor product with the oscillator vacuum gives the full vacuum.
An arbitrary lattice state is a finite linear combination of
\[
 a_1(-n_1)\cdots a_s(-n_s)e^\alpha,
 \qquad n_i>0,\quad \alpha\in L.
\]

The different summands have the same oscillator basis but different
Heisenberg zero-mode actions. On the ground states,
\begin{equation}
 a(n)e^\alpha=0\quad(n>0),\qquad
 a(0)e^\alpha=(a,\alpha)e^\alpha,
 \qquad a\in\mathfrak h.
 \label{eq:voa-momentum-ground-state}
\end{equation}
All zero modes commute with the creation modes, so every state in
this summand has the same momentum $\alpha$. Positive modes act
on oscillator excitations through the Heisenberg commutators, just
as in $M(1)$. These are the several-boson versions of the charged
Fock modules $F_p$ from Chapter~\ref{ch:cft}.

The free-boson stress tensor determines the conformal grading.
For an orthonormal basis $u_1,\ldots,u_r$ of $\mathfrak h$,
its zero mode is
\[
 L_0=\frac12\sum_{j=1}^r u_j(0)^2
       +\sum_{n>0}\sum_{j=1}^r u_j(-n)u_j(n).
\]
Only finitely many terms act nontrivially on any fixed oscillator
monomial.
On $e^\alpha$ the second sum vanishes, and the first gives
$\tfrac12\sum_j(u_j,\alpha)^2=\alpha^2/2$. Since a creation
mode $a(-n)$ adds weight $n$, a nonzero oscillator monomial has
\[
 \wt\bigl(a_1(-n_1)\cdots a_s(-n_s)e^\alpha\bigr)
 =\frac{\alpha^2}{2}+n_1+\cdots+n_s.
\]
Now the role of evenness is visible: every momentum ground state has
integer weight, so all conformal weights occurring in $V_L$ are integers.

The exponential notation is motivated by the associated field
$Y(e^\alpha,z)$, the lattice counterpart of the normal-ordered
free-boson exponential from Chapter~\ref{ch:cft}, customarily written
$:\!\exp(i\alpha\cdot X(z))\!:$ in physics notation.
Here $X=(X_1,\ldots,X_r)$ is a collection of independent chiral
bosons with $X_j(z)X_k(w)\sim-\delta_{jk}\log(z-w)$, and
$\alpha\cdot X=\sum_j(u_j,\alpha)X_j$. In the lattice theory
this notation includes the momentum-shift and cocycle factors
described below. The state $e^\alpha$ is the input to the
state--field map; $Y(e^\alpha,z)$ is the resulting operator-valued
field, which takes momentum $\beta$ to momentum $\alpha+\beta$.

\begin{example}[Momentum and oscillator excitations in rank one]
\label{ex:voa-momentum-labels}
Let $L=\Z\alpha$ with $\alpha^2=2$. Its momentum labels are
$e^{m\alpha}$ for $m\in\Z$, with ground-state weights $m^2$.
The first examples are
\[
\begin{array}{c|c|c}
 \text{state}&\text{momentum}&\text{weight}\\ \hline
 \one&0&0\\
 e^\alpha,\ e^{-\alpha}&\alpha,\ -\alpha&1\\
 \alpha(-1)\one&0&1\\
 \alpha(-1)e^\alpha&\alpha&2.
\end{array}
\]
Thus there are three independent weight-one states: two charged
ground states and one neutral oscillator state. A superposition such
as $e^\alpha+e^{-\alpha}$ is a sum of two different momentum
states, whereas $e^{\alpha+(-\alpha)}=e^0$ labels the zero-momentum
ground state. Keeping this distinction in mind prevents ordinary
exponential manipulations from being applied to vector addition.
\end{example}

\subsection{Multiplying momentum labels and choosing a cocycle}

Adding charges suggests the ordinary group-algebra multiplication
$e^\alpha e^\beta=e^{\alpha+\beta}$. The lattice construction
allows a scalar factor in that rule:
\[
 e^\alpha e^\beta=\epsilon(\alpha,\beta)e^{\alpha+\beta},
 \qquad \epsilon:L\times L\longrightarrow\C^\times.
\]
Here $\C^\times$ denotes the nonzero complex numbers. We require
$\epsilon(0,\alpha)=\epsilon(\alpha,0)=1$, so $e^0$ is the
identity for this multiplication. The resulting algebra is called
a \emph{twisted group algebra}. The concrete choice below takes
values in $\{1,-1\}$.

This multiplication records how a charged field shifts the momentum
label. For example, multiplication by $e^\alpha$ sends the label
$e^\beta$ to $\epsilon(\alpha,\beta)e^{\alpha+\beta}$.
It supplies one ingredient of $Y(e^\alpha,z)$; oscillator operators
and powers of $z$ supply the others. Consequently this auxiliary
group-algebra product is different from the VOA product
$e^\alpha_{(-1)}e^\beta$, which is a particular coefficient of
the full field $Y(e^\alpha,z)e^\beta$.
Associativity of this multiplication requires the two-cocycle
identity
\[
 \epsilon(\alpha,\beta)\epsilon(\alpha+\beta,\gamma)
 =\epsilon(\beta,\gamma)\epsilon(\alpha,\beta+\gamma).
\]
The two sides are the coefficients of $e^{\alpha+\beta+\gamma}$
obtained by multiplying three basis labels in the two possible
orders of association. A normalized function satisfying this identity
is called a \emph{two-cocycle}.
Locality of the exponentials also requires
\begin{equation}
 \frac{\epsilon(\alpha,\beta)}{\epsilon(\beta,\alpha)}
 =(-1)^{(\alpha,\beta)}.
 \label{eq:voa-cocycle-commutator}
\end{equation}
This sign compensates for interchanging the factor
$(z-w)^{(\alpha,\beta)}$ in the Gaussian OPE. The cocycle is
therefore part of a consistent choice of field phases, not an
optional decoration.

In the simply laced root bases used later, with Gram matrix $G$ and
diagonal entries two, write $a=(a_i)$ and $b=(b_i)$ for the integer
coordinate columns of two lattice vectors. The entry $G_{ij}$ is
the pairing of the $i$th and $j$th basis vectors. One convenient
cocycle is
\begin{equation}
 \epsilon(a,b)=(-1)^{\sum_i a_i b_i+
                              \sum_{i>j}a_i b_jG_{ij}}.
 \label{eq:voa-root-cocycle}
\end{equation}
Its exponent is bilinear modulo two, so the cocycle identity follows
by expanding both sides. The difference between the exponents for
$(a,b)$ and $(b,a)$ is $a^{\mathsf t}Gb$ modulo two, proving
\eqref{eq:voa-cocycle-commutator}. Moreover,
$\epsilon(a,-a)=(-1)^{a^2/2}$ in this convention.

Rescaling the basis $e^\alpha\mapsto\lambda(\alpha)e^\alpha$,
with $\lambda:L\to\C^\times$ and $\lambda(0)=1$, changes
$\epsilon$ by the coboundary
$\lambda(\alpha)\lambda(\beta)/\lambda(\alpha+\beta)$.
The condition at zero keeps the vacuum fixed. A rescaling of unit
modulus is a rephasing. The resulting VOA is isomorphic. A calculation that compares two
tables of coefficients must first align these gauges.

\section{The exponential formula and a root calculation}

The full vertex operator contains creation oscillators, annihilation
oscillators, and the charge-shift operator introduced above.
On a pure momentum state $e^\beta$, the positive oscillators
vanish and the zero modes measure $\beta$, as in
\eqref{eq:voa-momentum-ground-state}.
Acting on a state without oscillator excitations, the exponential
field reduces to the particularly useful formula
\begin{equation}
 Y(e^\alpha,z)e^\beta
 =\epsilon(\alpha,\beta)z^{(\alpha,\beta)}
   \exp\!\left(\sum_{n>0}\frac{\alpha(-n)}n z^n\right)
   e^{\alpha+\beta}.
 \label{eq:voa-exponential-action}
\end{equation}
The prefactor $z^{(\alpha,\beta)}$ sets the first possible mode,
and the exponential then generates the oscillator corrections.
For an oscillator-excited input, the needed additional identity is
\[
 [a(m),Y(e^\alpha,z)]=(a,\alpha)z^mY(e^\alpha,z).
\]
For example, moving one creation mode past the field gives, for $r>0$,
\[
 Y(e^\alpha,z)b(-r)e^\beta
 =\bigl(b(-r)-(b,\alpha)z^{-r}\bigr)Y(e^\alpha,z)e^\beta.
\]
Repeated use of this identity and the Heisenberg relations reduces
an arbitrary Fock monomial to the displayed pure-momentum formula.
At a prescribed output weight, only finitely many terms of the
exponential can contribute. This supplies an exact finite procedure
for each mode product.

\begin{example}[The affine $\mathfrak{sl}_2$ inside $V_{A_1}$]
\label{ex:voa-a1}
Let $L=\Z\alpha$ with $\alpha^2=2$. In the cocycle convention
above, $\epsilon(\alpha,-\alpha)=-1$, so
\[
 Y(e^\alpha,z)e^{-\alpha}
 =-z^{-2}\left[\one+\alpha(-1)z+
 \frac{\alpha(-1)^2+\alpha(-2)}2z^2+\cdots\right].
\]
Consequently
\[
 e^\alpha_{(1)}e^{-\alpha}=-\one,
 \qquad e^\alpha_{(0)}e^{-\alpha}=-\alpha(-1)\one.
\]
Define the states $H=\alpha(-1)\one$, $E=e^\alpha$, and
$F=-e^{-\alpha}$, and abbreviate their fields by
$H(z)=Y(H,z)$, $E(z)=Y(E,z)$, and $F(z)=Y(F,z)$. Then
\begin{align*}
 H(z)H(w)&\sim\frac2{(z-w)^2},\\
 H(z)E(w)&\sim\frac{2E(w)}{z-w},&
 H(z)F(w)&\sim-\frac{2F(w)}{z-w},\\
 E(z)F(w)&\sim\frac1{(z-w)^2}+\frac{H(w)}{z-w}.
\end{align*}
These are the affine $\mathfrak{sl}_2$ relations at level one.
The minus sign in $F$ is forced by the chosen cocycle and ensures
$[E_0,F_0]=H_0$.
\end{example}

At weight one, a positive even lattice contributes its rank from
the oscillator states and one state for every lattice vector
of squared norm two. Therefore
\begin{equation}
 \dim(V_L)_1=\rank L+\#\{\alpha\in L:\alpha^2=2\}.
 \label{eq:voa-weight1-lattice}
\end{equation}
For $A_1$ this is $1+2=3$. For $E_8$ it is $8+240=248$.
This elementary count will make the properness of the six-current
subalgebra a physically visible statement.

\section{Modules, fusion, and intertwining fields}
\index{module!vertex algebra}\index{intertwining operator}

A representation of a Lie algebra tells us how its generators act on
another vector space. A representation of a vertex algebra does the
same for all field modes at once. A \emph{$V$-module} $M$ carries a map
\[
 Y_M:V\longrightarrow\Hom_{\C}(M,M((z))),\qquad
 v\longmapsto Y_M(v,z),
\]
with mode expansion
\[
 Y_M(v,z)=\sum_{n\in\Z}v^M_{(n)}z^{-n-1},\qquad v\in V,
\]
whose coefficients $v^M_{(n)}:M\to M$ act on the module.
These fields obey the vacuum identity
$Y_M(\one,z)=\mathrm{id}_M$, lower truncation, translation, and the
same Jacobi or OPE compatibility identities as the fields of $V$.
Here translation includes
$Y_M(\mathsf Tv,z)=\partial_zY_M(v,z)$, where $\mathsf T$
acts on the input $v\in V$.
In particular, the conformal vector acts by Virasoro operators on $M$.
The module need not contain a vacuum state. Nor does its definition
assign an endomorphism-valued field to every vector of $M$: the inputs
of $Y_M$ still belong to the vacuum algebra $V$.

A module is \emph{simple} if it has no nonzero proper subspace
invariant under all field modes. In the simple ordinary
positive-energy modules used in the examples, $L_0^M$ is diagonalizable
and
\[
 M=\bigoplus_{N\geq0}M_{h_M+N},\qquad
 M_\lambda=\{u\in M:L_0^Mu=\lambda u\},
\]
with finite-dimensional eigenspaces and a lowest nonzero one
$M_{h_M}$. This lowest space is the \emph{top}; a state in
$M_{h_M+N}$ has \emph{grade} $N$ above it. Thus a module's top
can contain several states even when the module is simple. The
underlying module axioms alone do not impose this energy decomposition;
the final section distinguishes the various grading assumptions.

A \emph{module homomorphism} $f:M\to N$ is a linear map commuting
with these actions:
\[
 fY_M(v,z)=Y_N(v,z)f\qquad(v\in V).
\]
This is the meaning of $\Hom_V(M,N)$. For simple ordinary modules
with finite-dimensional energy eigenspaces, Schur's lemma makes it
zero for inequivalent modules and one-dimensional for equivalent
modules. For the scalar-endomorphism conclusion, restrict an
endomorphism to a nonzero finite energy space and choose an eigenvalue;
the kernel after subtracting that scalar is a nonzero submodule,
so simplicity makes it the whole module.

A \emph{chiral intertwining operator} is a different kind of object.
Here we discuss ordinary modules, with diagonalizable conformal energy.
An intertwiner of type $\binom{P}{M\ N}$ assigns to $u\in M$ a field
\[
 \mathcal Y(u,z):N\longrightarrow P\{z\}.
\]
For fixed $u$, its coefficient operators map $N$ to $P$, and its
action on a vector of $N$ is a vector-valued formal series in
$P\{z\}$. This notation permits the complex powers of $z$
prescribed by the module weights; it does not denote evaluation
of an analytic function. Generalized modules may require logarithmic
intertwiners as well, which are outside these examples.
The assignment is linear in $u$ and satisfies lower truncation, the
translation identity
$\mathcal Y(L_{-1}u,z)=\partial_z\mathcal Y(u,z)$, and compatibility
with the $V$-actions on all three modules. In physical language, moving
a chiral-algebra field past this insertion must agree with first using
its OPE with the inserted state. Thus a charged field can change the
sector on which it acts, even though every ordinary $V$ field preserves
each $V$-module. Products of such charged fields supply the nonvacuum
conformal blocks used later in the book.

The energy commutator is a useful part of this definition:
\[
 L_0^P\mathcal Y(u,z)-\mathcal Y(u,z)L_0^N
 =\mathcal Y(L_0^Mu,z)+z\partial_z\mathcal Y(u,z).
\]
If $u,v$ have weights $h_u,h_v$, a term $z^r w$ in
$\mathcal Y(u,z)v$ must therefore satisfy
\begin{equation}
 r=h_w-h_u-h_v.
 \label{eq:voa-intertwiner-power}
\end{equation}
For top states in three modules, the possible top-to-top power is
$h_P-h_M-h_N$. Its coefficient can vanish for particular choices of
top components. The first nonzero term then comes from a higher grade
of the same intermediate module; it need not signal a different
fusion channel.

The \emph{contragredient module} $M'$ is the graded dual
$\bigoplus_\lambda M_\lambda^*$ equipped with the dual
vertex-algebra action. It uses the dual of each finite energy space,
not the full algebraic dual of an infinite direct sum.
The VOA is \emph{self-contragredient} when $V'\cong V$ as
$V$-modules. This is one of the hypotheses in the next theorem.

To identify these spaces with categorical fusion, we use the standard
rational VOA tensor theorem under its usual hypotheses: simplicity,
CFT type, self-contragredience, rationality, $C_2$-cofiniteness, and
positive lowest weights for nonvacuum simple modules. The finiteness
terms are explained at the end of this chapter. In this setting the
ordinary modules form a semisimple rigid tensor category. Denote its
fusion product by $\boxtimes_V$ to distinguish it from the ordinary
tensor product of vector spaces. For simple modules,
\begin{equation}
 M\boxtimes_VN\cong\bigoplus_P N_{MN}^{\ P}\,P,
 \qquad
 N_{MN}^{\ P}
 =\dim\mathcal I\binom{P}{M\ N}
 =\dim\Hom_V(M\boxtimes_VN,P),
 \label{eq:voa-intertwiner-fusion}
\end{equation}
where $\mathcal I\binom{P}{M\ N}$ is the space of chiral
intertwiners of the indicated type. These are fusion maps
$M\boxtimes_VN\to P$. Chapter~\ref{ch:fusion} uses the oppositely
oriented splitting space $\Hom(P,M\boxtimes_VN)$; composition
pairs the two spaces nondegenerately. In the unitary setting,
adjoints convert orthonormal fusion bases to splitting bases.
The infinite tower of descendant
states belongs inside each module $P$; the finite integer
$N_{MN}^{\ P}$ counts independent ways of coupling to that module.
It does not count the dimension of $M\otimes_\C N$.

The contragredient $M'$ supplies the categorical dual in this setting.
A bar over a factor
in a full-CFT Hilbert space instead indicates an antiholomorphic copy.
The chosen right-moving pairing relates these notions; they are not
the same notation for taking complex conjugates of coordinates.
All twelve Haagerup center labels are self-dual, which will simplify
the displayed diagonal pairings without eliminating the underlying
duality convention.

\section{Lattice modules and a rank-one example}
\index{discriminant group}\index{quantum dimension!versus top multiplicity}

The dual lattice is
$L^*=\{x\in L\otimes\R:(x,L)\subset\Z\}$.
For a positive even lattice, the simple modules are
$V_{\lambda+L}$, indexed by $\lambda\in L^*/L$. Their conformal
weight is the minimum of $x^2/2$ in the coset, and their twist is
$e^{\pi i\lambda^2}$. Fusion is addition of cosets:
\[
 V_{\lambda+L}\boxtimes V_{\mu+L}
 \cong V_{\lambda+\mu+L}.
\]
Every categorical dimension is one. The representation category is
therefore pointed, even when some module tops have dimension greater
than one because a coset has several shortest vectors.
Concretely, $V_{\lambda+L}$ has the same oscillator monomials as
$V_L$, but its momentum basis is labelled by $\lambda+L$.
A vacuum-algebra field of momentum $\alpha\in L$ sends momentum
$\beta\in\lambda+L$ to $\alpha+\beta$ in the same coset.
This is how a charged lattice module can be preserved by every
$V_L$ field while containing no vacuum state of its own.

\begin{example}[A discriminant group one can compute by hand]
Let $L=\Z\alpha$, $\alpha^2=2k$ with $k\in\Z_{>0}$. Then
$L^*=\frac1{2k}\Z\alpha$, so $L^*/L\cong\Z/(2k)$.
The module labelled by $m$ has weight modulo one
\[
 q(m)=\frac{m^2}{4k}\pmod1.
\]
Its actual lowest weight uses the representative of $m$ closest
to zero modulo $2k$. The finite quadratic function $q$ records
twists, while the shortest-vector problem determines the integer
part of the energy. The distinction will reappear for the Haagerup
module whose twist exponent is $2/13$ but whose lowest weight is
$15/13$.
\end{example}

The character of a lattice module is
\begin{equation}
 \chi_{\lambda+L}(\tau)
 =\frac{\Theta_{\lambda+L}(\tau)}{\eta(\tau)^r},
 \qquad
 \Theta_{\lambda+L}=\sum_{x\in\lambda+L}q^{x^2/2}.
 \label{eq:voa-lattice-character}
\end{equation}
The numerator counts momenta and the denominator counts oscillators.
Poisson summation produces the modular $S$ matrix
\[
 S_{\lambda\mu}=|L^*/L|^{-1/2}e^{-2\pi i(\lambda,\mu)}.
\]
Thus the finite discriminant form is a direct bridge between a
Euclidean lattice and a modular tensor category.

\begin{example}[Poisson summation and the two sectors of the root lattice]
\label{ex:voa-poisson}
For $L=\Z\alpha$, $\alpha^2=2k$, write the theta series as
\[
 \Theta_m(\tau)=\sum_{\ell\in\Z}
 \exp\left(\frac{\pi i\tau}{2k}(2k\ell+m)^2\right),
 \qquad m\in\Z/(2k).
\]
Poisson summation replaces a sum of Gaussian samples on a lattice
by the Gaussian samples on its dual lattice. Applying the Gaussian
Fourier transform to the arithmetic progression $m+2k\Z$ gives
\[
 \Theta_m(-1/\tau)=\sqrt{\frac{-i\tau}{2k}}
 \sum_{r\in\Z}e^{\pi i rm/k}e^{\pi i\tau r^2/(2k)}.
\]
Here $\operatorname{Im}\tau>0$, and the analytic branch of
$\sqrt{-i\tau}$ is chosen positive at $\tau=it$ for $t>0$.
Positive definiteness of the lattice makes these theta sums
convergent in this domain.
Split the integer $r$ into its residue modulo $2k$ and an
integer multiple of $2k$. Reindex the residue by its negative;
$\Theta_{-n}=\Theta_n$ because reversing every momentum does
not change its norm. The result is
\begin{equation}
 \Theta_m(-1/\tau)=
 \sqrt{\frac{-i\tau}{2k}}
 \sum_{n=0}^{2k-1}e^{-\pi i mn/k}\Theta_n(\tau).
 \label{eq:voa-rankone-poisson}
\end{equation}
One can see both factors directly. The Fourier transform of the
Gaussian supplies $\sqrt{-i\tau}$; the spacing $2k$ leaves the
factor $(2k)^{-1/2}$. Splitting the dual sum into residue classes
modulo $2k$ leaves the finite Fourier phase. The sign of that phase
is the convention used here for $S$; reindexing $n\mapsto-n$
does not change these unflavoured theta series. Dividing by
$\eta(-1/\tau)=\sqrt{-i\tau}\eta(\tau)$ cancels the
continuous Gaussian factor. The modular transform of the character
is therefore an ordinary finite Fourier transform.

At $k=1$, there are two sectors, $0$ and $1$, and
\[
 S=\frac1{\sqrt2}\begin{pmatrix}1&1\\1&-1\end{pmatrix},
 \qquad
 T=e^{-\pi i/12}\begin{pmatrix}1&0\\0&i\end{pmatrix}.
\]
The common phase in $T$ comes from $c=1$; the relative phase
$i$ comes from the nontrivial lowest weight $h=1/4$.
Their characters begin
\begin{align*}
 \chi_0&=q^{-1/24}(1+3q+4q^2+\cdots),\\
 \chi_1&=q^{5/24}(2+2q+6q^2+\cdots).
\end{align*}
The two lowest states in the second sector are
$e^{\alpha/2}$ and $e^{-\alpha/2}$. Nevertheless,
$S_{10}/S_{00}=1$, its quantum dimension. Its square fuses to
the vacuum sector because adding its lattice coset to itself
gives $L$.

As a further check, the normalized twist sum is
$(1+i)/\sqrt2=e^{2\pi i/8}$, the phase appropriate to central
charge one. Chapter~\ref{ch:modular} explains this relation in
general. A nontrivial twist therefore need not mean a nontrivial
quantum dimension, just as a nontrivial top multiplicity need not.
\end{example}

\begin{example}[One fusion channel, two different leading powers]
\label{ex:voa-half-charge}
Let $L=\Z\alpha$ with $\alpha^2=2$, and put $V=V_L$.
The nontrivial simple module is
$M_1=V_{L+\alpha/2}$. Its top has basis
\[
 v_+=e^{\alpha/2},\qquad v_-=e^{-\alpha/2},
 \qquad h_{M_1}=\frac14.
\]
Addition of lattice cosets gives
$M_1\boxtimes_VM_1=V$, so there is exactly one vacuum-module
intertwiner up to scale. This does not mean that the OPE of every
pair of top components begins with the vacuum \emph{state}.

Choose an intertwiner $\mathcal Y$ of type
$\binom{V}{M_1\ M_1}$ and a branch for its fractional powers.
The lattice exponential formula gives
\begin{align}
 \mathcal Y(v_+,z)v_-
 &=\kappa_{+-}z^{-1/2}
   \left(\one+\frac12\alpha(-1)\one\,z+O(z^2)\right),
   \label{eq:voa-half-charge-neutral}\\
 \mathcal Y(v_+,z)v_+
 &=\kappa_{++}z^{1/2}
   \left(e^\alpha+\frac12\alpha(-1)e^\alpha\,z+O(z^2)\right).
   \label{eq:voa-half-charge-charged}
\end{align}
The nonzero phase coefficients $\kappa_{+-},\kappa_{++}$ depend on
the chosen intertwiner, basis, and cocycle extension convention; they
are not fixed merely by the cocycle on the original lattice $L$.
With the usual unit exponential normalizations the common scale can
be chosen so these coefficients have unit modulus. Their values are
irrelevant to the leading-power distinction.

In the first product the charges add to zero and
$(\alpha/2,-\alpha/2)=-1/2$. Its vacuum coefficient is nonzero.
In the second product the charges add to $\alpha$ and
$(\alpha/2,\alpha/2)=1/2$. Charge conservation forbids a vacuum-state
coefficient. The first output is instead the weight-one current
$e^\alpha$ in the same vacuum module $V$. Equation~
\eqref{eq:voa-intertwiner-power} verifies both exponents:
\[
 0-\frac14-\frac14=-\frac12,
 \qquad
 1-\frac14-\frac14=\frac12.
\]
The gap of one is an intermediate-module grade. It does not change
the fusion rule or create a second independent intertwiner.

For comparison, the inequivalent modules have no nonzero
module homomorphism between them, whereas their charged field exists:
\[
 \Hom_V(M_1,V)=0,\qquad
 \dim\mathcal I\binom{V}{M_1\ M_1}=1.
\]
A mode-preserving linear map between modules and a charged field
joining three modules answer different representation-theoretic questions.
\end{example}

\section{Affine currents, Sugawara stresses, and commutants}
\index{Sugawara construction}\index{commutant}

For a simple Lie algebra $\mathfrak g$ with invariant form normalized
so long roots have squared length two, choose a basis $t^a$.
Its associated current states $J^a$ have weight one; write
$J^a(z)=Y(J^a,z)$. The \emph{level} $k$ is the coefficient
multiplying the invariant form in their double pole:
\[
 J^a(z)J^b(w)\sim\frac{k\kappa^{ab}}{(z-w)^2}
                   +\frac{f^{ab}{}_cJ^c(w)}{z-w}.
\]
The structure constants and metric
are defined by $[t^a,t^b]=f^{ab}{}_c t^c$ and
$\kappa^{ab}=(t^a,t^b)$. The inverse metric $\kappa_{ab}$ defines
the dual basis $t_a=\kappa_{ab}t^b$ and the currents
$J_a=\kappa_{ab}J^b$; repeated indices are summed.
At positive integer level $k$, write $L_k(\mathfrak g)$ for the
simple integrable affine VOA. It is the simple quotient of the
vacuum representation generated by these currents. Integrability
requires the raising and lowering operators for affine simple roots
to act locally nilpotently: repeated application to any fixed state
eventually gives zero. For $\mathfrak{sl}_2$, one resulting relation is
$(E_{(-1)})^{k+1}\one=0$. At level one this is precisely the
relation $E_{(-1)}E=0$ visible in the root-lattice construction.
Such relations are not imposed in the universal affine vacuum algebra.
The Sugawara stress of the simple quotient is
\[
 T_{\mathfrak g,k}=\frac1{2(k+h^\vee)}
       \sum_a:J^aJ_a:,
 \qquad c_{\mathfrak g,k}=\frac{k\dim\mathfrak g}{k+h^\vee}.
\]
The dual Coxeter number $h^\vee$ is a constant of the simple Lie
algebra: in this normalization the quadratic Casimir
$\sum_a t^at_a$ acts on the adjoint representation by $2h^\vee$.
It supplies the correction from the simple poles when the current
bilinears are contracted. For $\mathfrak{sl}_2$, $h^\vee=2$ and
$c=3k/(k+2)$.

If an affine subalgebra $A$ is embedded in $V$, its Sugawara stress
need not equal the ambient stress. Suppose the currents are ambient
weight-one primaries, their level is positive integral, and the
embedding respects the ambient unitary grading. Under these
coset-embedding hypotheses, the difference
$\omega_C=\omega_V-\omega_A$ is a commuting stress of charge
$c_V-c_A$. The \emph{commutant}
\begin{equation}
 \Com(A,V)=\{v\in V:a_{(n)}v=0\text{ for all }a\in A,n\ge0\}
 \label{eq:voa-commutant}
\end{equation}
is the algebra of fields whose OPE with $A$ has no singular part.
Equivalently one can test a suitable generating set, with the
appropriate closure justification.

This is a chiral construction. A commuting pair of stresses does
not by itself imply a tensor-product decomposition of $V$.
The ambient states can form a nontrivial extension of the tensor
product, and can branch with multiplicities. Nor does a commutant
automatically supply a modular-invariant full theory. Those are
separate questions about representations and sewing.

An important example is the Cartan current in
$L_k(\mathfrak{sl}_2)$, which generates a rank-one Heisenberg
subalgebra $M(1)$. Its commutant is called the \emph{parafermion VOA}:
\[
 K(\mathfrak{sl}_2,k)=\Com(M(1),L_k(\mathfrak{sl}_2)).
\]
The Heisenberg stress has charge one, so this
example supplies the subtraction requested in
Exercise~\ref{ex:voa-coset}. Its states must satisfy the full
commutant condition, not merely have zero Cartan charge.

\begin{example}[The diagonal affine coset]
Embed $\mathfrak{sl}_2$ diagonally in level $k$ and level one.
The diagonal currents have level $k+1$, since their double poles
add. The coset central charge is
\[
 \frac{3k}{k+2}+1-\frac{3(k+1)}{k+3}
 =1-\frac6{(k+2)(k+3)}.
\]
At $k=1$ this is $1/2$, the Ising central charge. The equality of
central charges alone would not identify the coset, but the full
Goddard--Kent--Olive (GKO) construction identifies this family with the unitary Virasoro
minimal models. At $k=39$ the charge is $286/287$, a value that
will appear in the Haagerup stress analysis.
\end{example}

\begin{example}[The three Ising weights and a branching grade]
\label{ex:voa-Ising-branching}
Let $L_k(\lambda)$ denote the integrable affine
$\mathfrak{sl}_2$ module of level $k$ and highest weight
$\lambda\in\{0,\ldots,k\}$. Its top has spin
$j=\lambda/2$. In the invariant-form normalization above, the
quadratic Casimir acts by $2j(j+1)$. Dividing by $2(k+2)$
in the Sugawara formula gives
\[
 h_k(\lambda)=\frac{j(j+1)}{k+2}
             =\frac{\lambda(\lambda+2)}{4(k+2)}.
\]
Under the diagonal level-two algebra, write
\[
 L_1(\lambda)\otimes L_1(\mu)
 =\bigoplus_{\nu=0}^2 L_2(\nu)\otimes B_{\lambda,\mu}^{\nu}.
\]
The spaces $B$ are coset modules. If their first diagonal
highest vector occurs at numerator grade $N$, its coset weight is
\begin{equation}
 h_C=h_1(\lambda)+h_1(\mu)+N-h_2(\nu).
 \label{eq:voa-coset-grade}
\end{equation}
Here $N$ is an integer excitation above the numerator's own
lowest energy, and cannot be dropped from the formula.

Only labels with $\lambda+\mu-\nu$ even are admissible:
the diagonal $SU(2)$ center acts on the numerator by
$(-1)^{\lambda+\mu}$ and on the denominator by $(-1)^\nu$.
Affine current descendants preserve this parity. The diagonal
GKO branching theorem supplies sufficiency and completeness in
this example, and identifies
\[
 (\lambda,\mu;\nu)\sim(1-\lambda,1-\mu;2-\nu).
\]
This is an identification of coset modules, not of numerator
modules. There are six admissible triples and three pairs, with
no fixed triple. Representatives and their weights are
\begin{center}
\begin{tabular}{c|c|c|c}
$(\lambda,\mu;\nu)$ & $N$ & $h_C$ & Ising sector\\
\hline
$(0,0;0)$ & $0$ & $0$ & $\one$\\
$(1,0;1)$ & $0$ & $\frac14-\frac3{16}=\frac1{16}$ & $\sigma$\\
$(1,1;0)$ & $0$ & $\frac14+\frac14=\frac12$ & $\epsilon$
\end{tabular}
\end{center}
The last two grade-zero branches are visible in the finite
$\mathfrak{sl}_2$ products: doublet times singlet is a doublet,
and doublet times doublet contains a singlet.

The equivalent label $(0,0;2)$ reveals the grade correction.
Let $J_1^a,J_2^a$ be the two level-one currents and $\Omega$
the product vacuum. The three nonzero states
\[
 w^a=(J_{1,(-1)}^a-J_{2,(-1)}^a)\Omega
\]
occur at numerator grade one. For $J_D^b=J_1^b+J_2^b$,
the two equal central terms cancel, giving
\[
 J_{D,(n)}^b w^a=0\quad(n>0),\qquad
 J_{D,(0)}^b w^a=f^{ba}{}_c w^c.
\]
Thus the $w^a$ form a diagonal affine spin-one top; its raising
component is a highest vector with $\nu=2$. There is no such
top at grade zero, whose space is only the vacuum. Consequently
\[
 h_C=0+0+1-h_2(2)=1-\frac12=\frac12,
\]
as required by $(0,0;2)\sim(1,1;0)$. The states $w^a$ carry
diagonal charge; the coset sector is their multiplicity space.
They are not themselves fields of the vacuum commutant.

The GKO theorem identifies these branching modules with the
three irreducible Ising Virasoro modules and the commutant with
the Ising chiral algebra. Central charge subtraction alone had
not established that statement: it contained no branching
selection, identification, or lowest-grade information.
\end{example}

\section{Positive forms and a change of gauge}
\index{PCT involution}\index{gauge!transported positive form}

A lattice VOA has a canonical positive Hermitian form and an
antilinear PCT involution $\theta$. PCT denotes the conjugation
entering the radial-adjoint identity; it preserves the vacuum and
conformal vector. The invariant
bilinear form often used in purely algebraic VOA calculations is
a different object. Taking complex conjugates of coordinates is
not, by itself, the PCT operation in an arbitrary basis.

An automorphism of a VOA is an invertible linear map $g:V\to V$
with
\[
 g\one=\one,\qquad g\omega=\omega,\qquad
 gY(a,z)g^{-1}=Y(ga,z).
\]
In a lattice VOA, scaling $e^\alpha$ by
$\lambda(\alpha)\in\C^\times$ and fixing the oscillators gives
such an automorphism when $\lambda$ is a lattice character:
$\lambda(\alpha+\beta)=\lambda(\alpha)\lambda(\beta)$.
This preserves the chosen cocycle. The earlier arbitrary rescaling
instead changes the cocycle presentation and supplies an isomorphism
between the corresponding presentations.

Suppose such a VOA automorphism $g$ changes the normalization of
the exponential states. Transport the positive form and involution
together:
\begin{equation}
 \langle v,w\rangle_g=\langle g^{-1}v,g^{-1}w\rangle,
 \qquad \theta_g=g\theta g^{-1}.
 \label{eq:voa-gauge-transport}
\end{equation}
The new form is positive because $g$ is invertible. It may make a
displayed basis nonorthonormal with unequal diagonal norms. If
$v=\sum a_jv_j$ with orthogonal $v_j$, then
$\norm v_g^2=\sum|a_j|^2\norm{v_j}_g^2$, not necessarily
$\sum|a_j|^2$.

A subalgebra stable under the transported PCT inherits vacuum
unitarity. One must check stability; it is not automatic for a
complex-generated subalgebra. Even after vacuum unitarity is known,
positivity on all modules and on fusion spaces remains a further
problem. The complete-unitarity chapter develops this hierarchy.

\section{Three finiteness notions that will recur}
\index{ordinary module}\index{regularity}\index{C2-cofiniteness@$C_2$-cofiniteness}

An \emph{ordinary module} has diagonalizable $L_0$, finite-dimensional
energy eigenspaces, and lower truncation in every integer-shift
class: for each weight $\lambda$, the spaces of weights
$\lambda-n$ vanish for sufficiently large integers $n$.
A simple ordinary module in the positive-energy setting used here
has spectrum contained in $h+\Z_{\ge0}$ for one lowest weight $h$.
Finite direct sums may involve several such translates; for example,
a direct sum of an integral-weight module and a module of lowest weight
$1/3$ has both integer and $1/3$-shifted energies.
A weak module need not have these grading
properties. A generalized module can contain nontrivial Jordan
blocks for $L_0$, producing logarithmic behavior in correlation
functions.

An \emph{admissible module} is a weak module with a compatible
nonnegative integer grading $M=\bigoplus_{n\ge0}M(n)$:
for homogeneous $a$,
\[
 a_{(m)}M(n)\subset M(n+\wt(a)-m-1),
 \qquad M(j)=0\quad(j<0).
\]
This grading need not be the $L_0$ eigenspace grading, and its pieces
are not required to be finite dimensional.
\index{admissible module}

The $C_2$ subspace and quotient are
\begin{equation}
 C_2(V)=\Span\{a_{(-2)}b:a,b\in V\},\qquad
 R_V=V/C_2(V).
 \label{eq:voa-c2}
\end{equation}
Here $[a]$ denotes the coset $a+C_2(V)$. The quotient has a
commutative associative product
$[a][b]=[a_{(-1)}b]$ and a Poisson bracket
$\{[a],[b]\}=[a_{(0)}b]$. A Poisson bracket is a bilinear
antisymmetric operation satisfying the Jacobi identity and a
derivation rule for the commutative product. The vertex-algebra
identities make these operations independent of the representatives
$a,b$. Derivatives vanish in the quotient because
$\mathsf Ta=a_{(-2)}\one$. The condition
$\dim R_V<\infty$ is called $C_2$-cofiniteness. It controls a
large amount of representation theory, but by itself does not
exclude logarithmic nonsemisimple modules.

Rationality asserts complete reducibility of admissible modules.
Regularity is the stronger-looking
statement that every weak module is a direct sum of ordinary simple
modules. The standard theorem used here is precise: for a VOA of
CFT type, rationality together with $C_2$-cofiniteness implies
regularity. This is the Abe--Buhl--Dong implication. Thus the
strong rationality established in Chapter~\ref{ch:regularity}
supplies the weak-module decomposition needed in
Chapter~\ref{ch:unitarity}. We still distinguish the definitions,
since $C_2$-cofiniteness or a finite list of ordinary modules
alone would not supply that implication.

For the Haagerup algebra, the proof of finiteness is not a general
principle that every subalgebra of a lattice VOA is rational.
Heisenberg subalgebras already show why that would be false.
Instead, an explicit positive identity among conformal vectors
provides the needed control. Understanding why that identity works
is one of the main tasks of Part II.

\section{Exercises}

\begin{exercise}\label{ex:voa-weight}
For $a$ of weight three and $b$ of weight four, list the possible
nonzero singular products $a_{(n)}b$ in a CFT-type VOA and their
weights. Which coefficient is forced to be a scalar multiple of
the vacuum? Which coefficient is a candidate weight-two field?

\par\smallskip\noindent{\footnotesize\normalfont
\hyperref[sol:ex:voa-weight]{Solution on page~\pageref*{sol:ex:voa-weight}.}}\par
\end{exercise}

\begin{exercise}\label{ex:voa-cocycle}
Prove the cocycle and commutator identities for
\eqref{eq:voa-root-cocycle}. Then show directly that multiplying
$\epsilon$ by a coboundary leaves the commutator ratio unchanged.

\par\smallskip\noindent{\footnotesize\normalfont
\hyperref[sol:ex:voa-cocycle]{Solution on page~\pageref*{sol:ex:voa-cocycle}.}}\par
\end{exercise}

\begin{exercise}\label{ex:voa-rankone}
For $L=\Z\alpha$, $\alpha^2=2$, compute the first three coefficients
of $\Theta_L\prod_{n\ge1}(1-q^n)^{-1}$ and of the corresponding
series for $L+\alpha/2$. Explain why the second module has a
two-dimensional top but categorical dimension one.

\par\smallskip\noindent{\footnotesize\normalfont
\hyperref[sol:ex:voa-rankone]{Solution on page~\pageref*{sol:ex:voa-rankone}.}}\par
\end{exercise}

\begin{exercise}\label{ex:voa-heisenberg-c2}
Show that the rank-one Heisenberg VOA has $C_2$ quotient
$\C[x]$, with $x$ represented by $J_{-1}\one$. Conclude that
strong finite generation does not imply $C_2$-cofiniteness.

\par\smallskip\noindent{\footnotesize\normalfont
\hyperref[sol:ex:voa-heisenberg-c2]{Solution on page~\pageref*{sol:ex:voa-heisenberg-c2}.}}\par
\end{exercise}

\begin{exercise}\label{ex:voa-coset}
Derive the central charge of $K(\mathfrak{sl}_2,k)$ by removing
the Cartan Heisenberg stress from $L_k(\mathfrak{sl}_2)$.
Evaluate it at $k=3$ and $k=13$. Explain why this arithmetic alone
does not identify a pair of unspecified weight-two fields with the
two parafermion stresses.

\par\smallskip\noindent{\footnotesize\normalfont
\hyperref[sol:ex:voa-coset]{Solution on page~\pageref*{sol:ex:voa-coset}.}}\par
\end{exercise}

\begin{exercise}\label{ex:voa-gauge}
Suppose $g$ multiplies three orthonormal exponential states by
$1/2,1/\sqrt{40},1/\sqrt{40}$. Compute the transported squared norm of
$v=e^a+\frac{\sqrt{13}+3i\sqrt3}{20}e^b+
\frac{\sqrt{13}-3i\sqrt3}{20}e^c$. Explain why the answer differs
from the sum of the three squared coefficient magnitudes.

\par\smallskip\noindent{\footnotesize\normalfont
\hyperref[sol:ex:voa-gauge]{Solution on page~\pageref*{sol:ex:voa-gauge}.}}\par
\end{exercise}
